% SOURCE: Chinese PDF viewer pp.174--177, through the Dirac continuity equation.
\section{Relativistic Equations of Motion}
Difficulties of nonrelativistic quantum mechanics:
\begin{enumerate}[label=\arabic*)]
\item It has no relativistic covariance (momentum and energy are not treated symmetrically, nor are space and time).
\item Spin cannot be introduced self-consistently within the theory.
\item Probability is conserved, and there is no particle creation or annihilation.
\end{enumerate}
Passing from nonrelativistic to relativistic quantum mechanics is not simply the replacement
\[
E=\frac{\vect p^{\,2}}{2m}\quad\longrightarrow\quad E^2=\vect p^{\,2}c^2+m^2c^4.
\]
In the course of its development, a theory for relativistic many-particle systems---quantum field theory---was eventually established, with relativistic quantum mechanics representing an intermediate stage.

Relativistic quantum mechanics is an important tool for studying some quantum systems (for example, the hydrogen atom). More importantly, it provides a foundation for establishing quantum field theory.

\subsection*{1) Natural units}
There are three independent dimensions commonly used in physics: length $[L]$, time $[T]$, and energy $[E]$. In relativistic quantum mechanics, especially in the quantum field theory developed from it, a new system of units is introduced. Take the speed of light $c$ as the unit of velocity and $\hbar$ as the unit of action. Then only one independent unit remains, which may be chosen to be energy $[E]$.

Thus in the new unit system one may set $\hbar=c=1$, so that $\hbar$ and $c$ no longer appear explicitly in calculations.

\subsection*{2) Klein--Gordon equation}
Consider the wave equation for a free particle.

Nonrelativistic energy:
\[
E=\frac{\vect p^{\,2}}{2m}.
\]
Relativistic energy:
\[
E^2=\vect p^{\,2}+m^2.
\]
Apply the canonical quantization prescription
\[
E\longrightarrow\ii\frac{\partial}{\partial t},\qquad
\vect p\longrightarrow-\ii\vect\nabla.
\]
One obtains the nonrelativistic Schr\"odinger equation
\[
\ii\frac{\partial\psi}{\partial t}=-\frac{\vect\nabla^{2}}{2m}\psi,
\]
and the relativistic Klein--Gordon equation
\[
(\partial^\mu\partial_\mu+m^2)\psi=0,\qquad
\partial^\mu\partial_\mu=\partial_t^2-\vect\nabla^2.
\]

Multiply the Klein--Gordon equation on the left by $\psi^*$, and multiply its complex conjugate on the right by $\psi$:
\[
\psi^*(\partial^\mu\partial_\mu+m^2)\psi=0,\qquad
\bigl[(\partial^\mu\partial_\mu+m^2)\psi^*\bigr]\psi=0.
\]
Subtracting these two equations gives
\begin{align*}
\psi^*(\partial^\mu\partial_\mu\psi)-(\partial^\mu\partial_\mu\psi^*)\psi&=0,\\
\partial^\mu\left[\psi^*(\partial_\mu\psi)-(\partial_\mu\psi^*)\psi\right]&=0.
\end{align*}
This is a conserved-current (continuity) equation:
\[
\partial^\mu j_\mu=0,\qquad
j_\mu=\psi^*(\partial_\mu\psi)-(\partial_\mu\psi^*)\psi.
\]
Alternatively,
\[
\partial_t\rho+\vect\nabla\cdot\vect j=0,
\]
where
\[
\rho=\frac{\ii}{2m}\left(\psi^*\frac{\partial\psi}{\partial t}-
\frac{\partial\psi^*}{\partial t}\psi\right),\qquad
\vect j=\frac{1}{2\ii m}
\left(\psi^*\vect\nabla\psi-\psi\vect\nabla\psi^*\right).
\]
However, $\rho$ is not positive definite, so it cannot be guaranteed that $\rho\geq0$. This is because the Klein--Gordon equation is second order in time. Both $\psi$ and $\partial\psi/\partial t$ are independent initial conditions; consequently, the above $\rho$ may be negative. This is the problem of a non-positive-definite probability in the Klein--Gordon equation. A way to solve the negative-probability difficulty is to treat the Klein--Gordon equation as a field equation (quantum field theory). Then the current is an electric current rather than a probability current, and the charge density can be negative. By contrast, the Schr\"odinger equation is first order in time, and $\rho=\psi^*\psi$ is positive definite.

Furthermore, a plane wave of arbitrary energy $E$,
\[
\psi=e^{-\ii p^\mu x_\mu}=e^{-\ii(Et-\vect p\cdot\vect x)},
\]
is a solution of the Klein--Gordon equation provided that $E^2=\vect p^{\,2}+m^2$. Hence plane waves with both
\[
E=+\sqrt{\vect p^{\,2}+m^2}\qquad\text{and}\qquad
E=-\sqrt{\vect p^{\,2}+m^2}
\]
solve the equation. This is the negative-energy difficulty of the Klein--Gordon equation. In classical mechanics a particle initially in a positive-energy state has no transitions and therefore no negative-energy difficulty. In quantum mechanics, transitions can lead to the negative-energy difficulty. Note that the Schr\"odinger equation has no negative-energy difficulty. A way to solve this problem is to introduce antiparticles (quantum field theory).

\subsection*{3) Dirac equation}
To address the negative-probability difficulty arising from a second-order time derivative, consider an equation that is first order in time. Relativistic covariance then requires treating the spatial derivatives as first order as well. Dirac therefore proposed the equation of motion for the wavefunction
\[
(\ii\gamma^\mu\partial_\mu-m)\psi=0.
\]
Here $\psi$ is a four-component wavefunction in Dirac space (as shown below, the four components correspond to positive and negative energy and the two spin directions); it is called a \emph{Dirac spinor wavefunction}. The $\gamma^\mu$ form a Lorentz four-vector of $4\times4$ matrices acting on Dirac space:
\[
\gamma^0=\begin{pmatrix}I&0\\0&-I\end{pmatrix},\qquad
\gamma^i=\begin{pmatrix}0&\sigma_i\\-\sigma_i&0\end{pmatrix}\quad (i=1,2,3).
\]
They satisfy
\begin{align*}
\{\gamma_\mu,\gamma_\nu\}&=\gamma_\mu\gamma_\nu+\gamma_\nu\gamma_\mu=2g_{\mu\nu}I,\\
(\gamma^\mu)^\dagger&=\gamma^0\gamma^\mu\gamma^0,\\
\gamma^\mu\gamma_\mu&=g^\mu{}_{\mu}I=4I,\\
\gamma^\mu\gamma^\nu\gamma_\mu&=-2\gamma^\nu,\\
\gamma^\mu\gamma^\nu\gamma^\rho\gamma_\mu&=4g^{\nu\rho}I,\\
\gamma^\mu\gamma^\nu\gamma^\rho\gamma^\sigma\gamma_\mu&=-2\gamma^\sigma\gamma^\rho\gamma^\nu.
\end{align*}
$g_{\mu\nu}$ is the metric tensor of four-dimensional Minkowski space:
\[
g_{\mu\nu}=g^{\mu\nu}=\operatorname{diag}(1,-1,-1,-1),\qquad
 g^\mu{}_{\nu}=g^{\mu\sigma}g_{\sigma\nu}=\operatorname{diag}(1,1,1,1).
\]
Since
\begin{align*}
(\ii\gamma^\mu\partial_\mu+m)(\ii\gamma^\nu\partial_\nu-m)\psi&=0,\\
-\bigl(\gamma^\mu\gamma^\nu\partial_\mu\partial_\nu+m^2\bigr)\psi&=0,\\
\left[-\frac12(\gamma^\mu\gamma^\nu+\gamma^\nu\gamma^\mu)\partial_\mu\partial_\nu-m^2\right]\psi&=0,\\
\left[-\frac12(2g^{\mu\nu})\partial_\mu\partial_\nu-m^2\right]\psi&=0,\\
(\partial^\mu\partial_\mu+m^2)\psi&=0,
\end{align*}
each component of the four-vector $\psi$ satisfies the Klein--Gordon equation. However, since the spinor also satisfies the Dirac equation, additional conditions apply. Thus positive- and negative-energy plane-wave solutions satisfying $E^2=\vect p^{\,2}+m^2$ are both solutions of the Dirac equation, showing that the Dirac equation also has the negative-energy difficulty.

Take the Hermitian conjugate of the Dirac equation:
\[
\psi^\dagger\bigl(\ii\gamma^{\mu\dagger}\overleftarrow{\partial}_\mu+m\bigr)=0.
\]
Using $(\gamma^\mu)^\dagger=\gamma^0\gamma^\mu\gamma^0$, $(\gamma^0)^2=I$, and defining the Dirac adjoint
\[
\overline\psi=\psi^\dagger\gamma^0,
\]
one obtains
\[
\overline\psi\bigl(\ii\gamma^\mu\overleftarrow{\partial}_\mu+m\bigr)=0.
\]
Multiplying the last equation on the right by $\psi$ and the Dirac equation on the left by $\overline\psi$, and adding them gives
\begin{align*}
\overline\psi\gamma^\mu(\overleftarrow\partial_\mu+\partial_\mu)\psi&=0,\\
\partial_\mu(\overline\psi\gamma^\mu\psi)&=0.
\end{align*}
Define the current
\[
j^\mu=\overline\psi\gamma^\mu\psi.
\]
The continuity equation becomes
\[
\partial_\mu j^\mu=0,\qquad
j^0=\rho=\overline\psi\gamma^0\psi=\psi^\dagger\psi\geq0
\quad\text{(positive definite)},\qquad
\vect j=\overline\psi\vect\gamma\psi.
\]
This shows that the Dirac equation does not suffer from a non-positive-definite probability.